% AA-54 — wyjątkowe binarne grupy wielościenne
\chapter{Niezmienniki binarnych grup wielościennych}
\label{aa:ch:exceptional-invariants}

W rozdziale~\ref{aa:ch:finite-invariants} wyznaczyliśmy
pierścienie niezmienników grup cyklicznych i dwucyklicznych.
Pozostały trzy grupy wyjątkowe z
Twierdzenia~\ref{aa:tw:binary-classification}.
Pracujemy w modelach kwaternionowych z
Przykładu~\ref{aa:prz:binary-quaternion-generators}.
Jeżeli $G'=LGL^{-1}$ dla $L\in\operatorname{GL}_2(\mathbb C)$,
to $P\mapsto P\circ L^{-1}$ jest izomorfizmem
$\mathbb C[x,y]^G\to\mathbb C[x,y]^{G'}$:
dla $g'=LgL^{-1}$ mamy
$(P\circ L^{-1})(g'v)=P(gL^{-1}v)=P(L^{-1}v)$.
Odwzorowanie odwrotne jest dane przez złożenie z $L$.
Oba odwzorowania zachowują stopnie.

\section{Szeregi Hilberta}

\begin{twierdzenie}[Szeregi Hilberta grup wyjątkowych]
\label{aa:tw:exceptional-hilbert}
Niech $G$ będzie jedną z grup $2T,2O,2I$ działających
naturalnie na $\mathbb C^2$. Dla pierścienia
$\mathbb C[x,y]^G$ z gradacją według stopnia mamy
\[
\begin{aligned}
 H_{2T}(z)&=\frac{1+z^{12}}{(1-z^6)(1-z^8)},\\
 H_{2O}(z)&=\frac{1+z^{18}}{(1-z^8)(1-z^{12})},\\
 H_{2I}(z)&=\frac{1+z^{30}}{(1-z^{12})(1-z^{20})}.
\end{aligned}
\]
\end{twierdzenie}
\begin{proof}
Niech $\phi=(1+\sqrt5)/2$. Ślady elementów
$\operatorname{SU}(2)$ są rzeczywiste.
Z konstrukcji nakrycia w rozdziale 52 wynika,
że dwa podniesienia obrotu o kąt $\theta$
mają ślady $\pm2\cos(\theta/2)$.
Korzystając z rozkładów rzędów w
Przykładzie~\ref{aa:prz:binary-order-distributions}
i z dwóch klas cykli długości pięć w $A_5$,
otrzymujemy tabelę liczności. W kolumnach oznaczonych
$\pm a$ podajemy liczność dla \emph{każdego} znaku.
\[
\begin{array}{c|rrrrrr}
 &\pm2&0&\pm1&\pm\sqrt2&\pm\phi&\pm(\phi-1)\\ \hline
2T&1&6&8&0&0&0\\
2O&1&18&8&6&0&0\\
2I&1&30&20&0&12&12
\end{array}
\]
Istotnie ślady pochodzące od obrotów rzędów
$2,3,4$ to odpowiednio $0,\pm1,\pm\sqrt2$.
Dwie klasy obrotów rzędu pięć mają po dwanaście
elementów i kąty, z dokładnością do znaku,
$2\pi/5$ oraz $4\pi/5$; ich podniesienia mają ślady
$\pm\phi$ oraz $\pm(\phi-1)$.

We wzorze Moliena
(Twierdzenie~\ref{aa:tw:molien}) łączymy składniki
odpowiadające przeciwnym śladom. Po podstawieniu
$s=z^2$ oznaczmy
\[
\begin{aligned}
 A&=\frac{2(1+s)}{(1-s)^2},&
 B&=\frac1{1+s},\\
 C&=\frac{2(1+s)}{1+s+s^2},&
 D&=\frac{2(1+s)}{1+s^2},\\
 E&=\frac{2(1+s)(2+s+2s^2)}
 {1+s+s^2+s^3+s^4}.
\end{aligned}
\]
Ostatni składnik jest sumą dwóch par odpowiadających
śladom $\pm\phi,\pm(\phi-1)$:
dla pary $\pm a$ mianownik po połączeniu to
$(1+s)^2-a^2s$. Iloczyn dwóch takich mianowników
dla $a=\phi,\phi-1$ wynosi
$(1+s)^4-3s(1+s)^2+s^2=1+s+s^2+s^3+s^4$,
gdyż $\phi^2+(\phi-1)^2=3$ oraz
$\phi^2(\phi-1)^2=1$.
Po sprowadzeniu obu par do tego mianownika licznik
wynosi $2(1+s)(2(1+s)^2-3s)$, co daje $E$.
Zatem
\[
\begin{aligned}
24H_{2T}&=A+6B+8C,\\
48H_{2O}&=A+18B+8C+6D,\\
120H_{2I}&=A+30B+20C+12E.
\end{aligned}
\]
Po sprowadzeniu do wspólnego mianownika dostajemy
odpowiednio
\[
 \frac{1+s^6}{(1-s^3)(1-s^4)},\quad
 \frac{1+s^9}{(1-s^4)(1-s^6)},\quad
 \frac{1+s^{15}}{(1-s^6)(1-s^{10})}
\]
dla $H_{2T},H_{2O},H_{2I}$. Równości można sprawdzić
przez wymnożenie obu stron przez iloczyn mianowników;
używa się $1-s^k=(1-s)(1+s+\cdots+s^{k-1})$.
Powrót do $z$ daje tezę.
\end{proof}

\section{Kryterium kompletności generatorów}

\begin{lemat}[Kryterium prezentacji kwadratowej]
\label{aa:lem:exceptional-presentation}
Niech $G\le\operatorname{SL}_2(\mathbb C)$ będzie
skończona, a $a,b,c\in\mathbb C[x,y]^G$ niezerowymi
jednorodnymi niezmiennikami dodatnich stopni
$p,q,r$. Załóżmy, że $a,b$ są algebraicznie niezależne,
$c^2=P(a,b)$ dla $P\in\mathbb C[U,V]$,
a $P(U,V)$ nie jest kwadratem w $\mathbb C(U,V)$.
Jeżeli
\[
 H_G(z)=\frac{1+z^r}{(1-z^p)(1-z^q)},
\]
to homomorfizm
\[
 \mathbb C[U,V,W]/(W^2-P(U,V))
 \longrightarrow\mathbb C[x,y]^G,
 \quad (U,V,W)\longmapsto(a,b,c)
\]
jest izomorfizmem.
\end{lemat}
\begin{proof}
Dzielenie przez wielomian moniczny w $W$ daje
jednoznaczną postać $L(U,V)+WM(U,V)$ elementu
ilorazu. Gdy jej obraz jest zerowy i $M\ne0$,
niezależność algebraiczna daje $M(a,b)\ne0$.
Wówczas $c=-L(a,b)/M(a,b)$, zatem po podniesieniu
do kwadratu $P$ jest kwadratem w $\mathbb C(U,V)$,
sprzecznie z założeniem. Jeśli $M=0$, niezależność
$a,b$ daje $L=0$. Homomorfizm jest więc iniektywny.
Nadajmy zmiennym $U,V,W$ wagi $p,q,r$.
Ponieważ $P(a,b)=c^2$ jest jednorodny stopnia $2r$
i $a,b$ są algebraicznie niezależne, każda część
$P$ o wadze różnej od $2r$ jest zerowa.
Relacja jest więc jednorodna dla tych wag,
a homomorfizm zachowuje gradację.
Jego obraz ma bazę jednorodną $a^ib^j$ i $ca^ib^j$,
$i,j\ge0$, oraz szereg Hilberta równy podanemu
$H_G$. Ponieważ obraz jest podprzestrzenią
przestrzeni niezmienników w każdym stopniu,
równość skończonych wymiarów daje równość
tych przestrzeni w każdym stopniu.
\end{proof}

\begin{lemat}[Dwa kryteria algebraiczne]
\label{aa:lem:exceptional-jacobian}
Niech $a,b\in\mathbb C[x,y]$.
Jeżeli $J(a,b)=a_xb_y-a_yb_x\ne0$, to $a,b$ są
algebraicznie niezależne. Ponadto niezerowy
$P\in\mathbb C[U,V]$ nieparzystego stopnia względem
$V$ nie jest kwadratem w $\mathbb C(U,V)$.
\end{lemat}
\begin{proof}
Gdyby istniała niezerowa relacja $Q(a,b)=0$,
wybralibyśmy $Q$ o najmniejszym stopniu całkowitym.
Nie jest stała. Różniczkowanie względem $x,y$ daje
układ jednorodny z niewiadomymi
$Q_U(a,b),Q_V(a,b)$ i wyznacznikiem $J(a,b)$.
Nad ciałem $\mathbb C(x,y)$ układ jest odwracalny,
więc oba wyrażenia znikają. Co najmniej jedna
pochodna wielomianu $Q$ jest niezerowa
(charakterystyka zero) i ma mniejszy stopień,
co przeczy minimalności.

Dla drugiej części traktujemy $P$ jako wielomian
w $V$ nad ciałem $\mathbb C(U)$. Gdyby
$P=(R/S)^2$ dla niezerowych $R,S\in\mathbb C(U)[V]$,
porównanie stopni w równości $PS^2=R^2$
dałoby $\deg_V P+2\deg_V S=2\deg_V R$,
co jest niemożliwe.
\end{proof}

\section{Grupy czworościenna i ośmiościenna}

Połóżmy
\[
\begin{aligned}
 f&=xy(x^4-y^4),\\
 g&=x^8+14x^4y^4+y^8,\\
 t&=x^{12}-33x^8y^4-33x^4y^8+y^{12}.
\end{aligned}
\]

\begin{lemat}[Sprawdzenie niezmienniczości]
\label{aa:lem:exceptional-small-forms}
W modelach z Przykładu~\ref{aa:prz:binary-quaternion-generators}
wielomiany $f,g,t$ są niezmiennikami $2T$.
Dla generatora $r=(1+\mathbf i)/\sqrt2$ grupy $2O$
podstawienie przez $M(r)$ wysyła je odpowiednio
na $-f,g,-t$. Zachodzi $t^2=g^3-108f^4$.
\end{lemat}
\begin{proof}
Macierze $M(\mathbf i),M(\mathbf j)$ dają
podstawienia $(ix,-iy)$ i $(y,-x)$; każdy z trzech
wielomianów jest przez nie zachowywany.
Dla $h=(1+\mathbf i+\mathbf j+\mathbf k)/2$ mamy
\[
 X=\tfrac{1+i}{2}(x+y),\qquad
 Y=\tfrac{-1+i}{2}(x-y).
\]
Stąd $XY=-(x^2-y^2)/2$,
$X^4=-(x+y)^4/4$, $Y^4=-(x-y)^4/4$.
Różnica czwartych potęg daje
$f(X,Y)=xy(x^4-y^4)$. Dla $g$ korzystamy z
\[
 16g(X,Y)=(x+y)^8+14(x^2-y^2)^4+(x-y)^8
          =16g(x,y).
\]
Dla $t$ mamy
\[
 64t(X,Y)=-(x+y)^{12}-(x-y)^{12}
 +33(x^2-y^2)^4\big((x+y)^4+(x-y)^4\big)
 =64t(x,y).
\]
Ostatnie równości wynikają z rozwinięcia dwumianowego.
Podstawienie przez $M(r)$ jest diagonalne z wpisami
$e^{\pi i/4},e^{-\pi i/4}$, więc monom $x^ay^b$
jest mnożony przez $e^{(a-b)\pi i/4}$. Daje to
wskazane znaki. Sprawdzanie przez podstawienie $M(q)$
zamiast $M(q)^{-1}$ jest równoważne niezmienniczości
względem $q$.

Na koniec połóżmy $A=x^4,B=y^4$.
Wówczas $f^4=AB(A-B)^4$,
$g=A^2+14AB+B^2$, $t=(A+B)(A^2-34AB+B^2)$.
Rozwinięcie obu stron daje
$g^3-t^2=108AB(A-B)^4=108f^4$.
\end{proof}

\begin{twierdzenie}[Prezentacje dla dwóch pierwszych grup]
\label{aa:tw:exceptional-e6-e7}
W oznaczeniach Lematu~\ref{aa:lem:exceptional-small-forms}
zachodzą izomorfizmy
\[
\begin{aligned}
 \mathbb C[x,y]^{2T}
 &\simeq\mathbb C[U,V,W]/(W^2-V^3+108U^4),\\
 \mathbb C[x,y]^{2O}
 &\simeq\mathbb C[U,V,W]/(W^2-UV^3+108U^3),
\end{aligned}
\]
gdzie generatory przechodzą odpowiednio na $(f,g,t)$
oraz $(f^2,g,ft)$.
\end{twierdzenie}
\begin{proof}
Mamy $f_x(1,0)=0$, $f_y(1,0)=1$,
$g_x(1,0)=8$, $g_y(1,0)=0$, stąd
$J(f,g)(1,0)=-8$. Zatem $f,g$ są algebraicznie
niezależne na mocy Lematu~\ref{aa:lem:exceptional-jacobian}.
Wielomian $V^3-108U^4$ ma nieparzysty stopień względem
$V$ i nie jest kwadratem. Stopnie $f,g,t$ to $6,8,12$.
Z Twierdzenia~\ref{aa:tw:exceptional-hilbert} i Lematu
\ref{aa:lem:exceptional-presentation} wynika pierwsza
prezentacja.

Grupa $2O$ jest generowana przez $h,r$, więc
$f^2,g,ft$ są jej niezmiennikami. Ich stopnie
wynoszą $12,8,18$. Niezależność $f^2,g$ wynika
z niezależności $f,g$ (podstawienie $U\mapsto U^2$
jest iniektywne na pierścieniu wielomianów).
Ponadto $(ft)^2=f^2(g^3-108f^4)$, a
$UV^3-108U^3$ ma stopień trzy względem $V$.
Stosujemy ponownie to samo kryterium prezentacji,
teraz z szeregiem Hilberta grupy $2O$.
\end{proof}

\section{Grupa dwudziestościenna}

Dla modelu $2I$ z rozdziału 52 potrzebny jest
inny niezmiennik stopnia dwanaście.
Oznaczmy $\phi=(1+\sqrt5)/2$ i połóżmy
\[
 F=\sqrt5\,t-22f^2,\quad
 H=F_{xx}F_{yy}-F_{xy}^2,\quad K=J(F,H).
\]
Ostatnie dwa wielomiany mają stopnie odpowiednio
$20$ i $30$, o ile nie są zerowe.

\begin{lemat}[Niezmienniczość hesjanu i jakobianu]
\label{aa:lem:exceptional-covariants}
Niech $G\le\operatorname{SL}_2(\mathbb C)$,
a $P,Q\in\mathbb C[x,y]^G$. Wówczas
$\det D^2P$ oraz $J(P,Q)$ są niezmiennikami $G$.
\end{lemat}
\begin{proof}
Dla $L\in G$ równość $P(Lv)=P(v)$ i reguła
łańcuchowa dają
$D^2P(v)=L^{\mathsf T}D^2P(Lv)L$.
Po wzięciu wyznaczników i użyciu $\det L=1$
otrzymujemy niezmienniczość hesjanu. Analogicznie
$D(P,Q)(v)=D(P,Q)(Lv)L$, więc wyznacznik
macierzy pochodnych pierwszego rzędu też jest
niezmiennikiem.
\end{proof}

\begin{lemat}[Wielomian stopnia dwanaście]
\label{aa:lem:exceptional-icosahedral-form}
W modelu $2I$ z Przykładu~\ref{aa:prz:binary-quaternion-generators}
wielomian $F$ jest niezmiennikiem.
\end{lemat}
\begin{proof}
Niezmienniczość względem $h,\mathbf i$ wynika
z Lematu~\ref{aa:lem:exceptional-small-forms}.
Dla trzeciego generatora
$u=\phi/2+(\phi-1)\mathbf j/2+\mathbf k/2$
podstawienie przez $M(u)$ daje
\[
 X=\frac{\phi x+(\phi-1+i)y}{2},\quad
 Y=\frac{(1-\phi+i)x+\phi y}{2}.
\]
Podajemy dokładną procedurę sprawdzenia współczynników,
aby nie zakładać niezmienniczości z samego stopnia.
Niech $A,B,C,D$ będą kolejno współczynnikami
$x,y$ w $2X$ i $2Y$. Dla
$F=\sum_{k=0}^{12}a_kx^{12-k}y^k$ współczynnik
$x^{12-l}y^l$ w $2^{12}F(X,Y)$ wynosi
\[
 \sum_{k=0}^{12}a_k
 \sum_{\substack{j+v=l\\0\le j\le12-k\\0\le v\le k}}
 \binom{12-k}{j}\binom{k}{v}
 A^{12-k-j}B^jC^{k-v}D^v.
\]
Używając $\phi^2=\phi+1$, $i^2=-1$,
$\sqrt5=2\phi-1$, po podzieleniu przez $2^{12}$
otrzymujemy dla parzystych $l$ tabelę
\[
\begin{array}{c|rrrrrrr}
l&0&2&4&6&8&10&12\\ \hline
a_l&\sqrt5&-22&-33\sqrt5&44&-33\sqrt5&-22&\sqrt5
\end{array}
\]
oraz zero dla nieparzystych $l$.
Jest to dokładnie ciąg współczynników $F$.
Zatem $F(X,Y)=F(x,y)$, a wskazane trzy
kwaterniony generują $2I$.
\end{proof}

\begin{twierdzenie}[Prezentacja dla grupy dwudziestościennej]
\label{aa:tw:exceptional-e8}
Niech $F,H,K$ będą wielomianami określonymi
powyżej dla modelu $2I$ z
Przykładu~\ref{aa:prz:binary-quaternion-generators}.
Wówczas
\[
 \mathbb C[x,y]^{2I}\simeq
 \frac{\mathbb C[U,V,W]}
 { (W^2+129534197760U^5+\frac{400}{121}V^3) },
 \qquad (U,V,W)\longmapsto(F,H,K).
\]
\end{twierdzenie}
\begin{proof}
Lematy~\ref{aa:lem:exceptional-covariants}
i~\ref{aa:lem:exceptional-icosahedral-form}
zapewniają niezmienniczość trzech wielomianów.
Relację wyprowadzimy z wymiaru części stopnia $60$
oraz dwóch współczynników. Z definicji pochodnych mamy
\[
\begin{array}{c|rr}
 & [x^d] & [x^{d-2}y^2] \\ \hline
 F,\ d=12 & \sqrt5 & -22\\
 H,\ d=20 & -5808\sqrt5 & -367840
\end{array}
\]
oraz $[x^{30}]K=0$,
$[x^{29}y]K=-13939200\sqrt5$.
Dokładniej, dla pierwszych trzech parzystych
współczynników $F$, oznaczonych $a_0,a_2,a_4$,
mamy
\[
 [x^{20}]H=264a_0a_2,\qquad
 [x^{18}y^2]H=1584a_0a_4-220a_2^2.
\]
Wynika to z wymnożenia pochodnych drugiego rzędu.
Dla $a_0=\sqrt5$, $a_2=-22$, $a_4=-33\sqrt5$
otrzymujemy liczby z tabeli.
Współczynnik przy $x^{29}y$ w $K$ wynosi
$12\sqrt5\cdot2(-367840)-2(-22)\cdot20(-5808\sqrt5)$.
Zatem $K=J(F,H)\ne0$ i $F,H$ są algebraicznie niezależne
na mocy Lematu~\ref{aa:lem:exceptional-jacobian}.

Z szeregu Hilberta $2I$ wynika, że część stopnia $60$
ma wymiar dwa: jedyne rozwiązania
$12i+20j+30e=60$, $i,j\ge0$, $e=0,1$,
to $(5,0,0)$ i $(0,3,0)$.
Wielomiany $F^5,H^3$ są liniowo niezależne
i stanowią więc jej bazę. Istnieją $\alpha,\beta$
takie, że $K^2=\alpha F^5+\beta H^3$.
Porównanie współczynników $x^{60}$ oraz $x^{58}y^2$
daje układ
\[
\begin{aligned}
0&=\alpha(\sqrt5)^5+\beta(-5808\sqrt5)^3,\\
5(13939200)^2
 &=5\alpha(\sqrt5)^4(-22)
   +3\beta(-5808\sqrt5)^2(-367840).
\end{aligned}
\]
Rozwiązanie to $\alpha=-129534197760$, $\beta=-400/121$
(oba niezerowe). Wielomian $\alpha U^5+\beta V^3$
ma nieparzysty stopień względem $V$, więc nie jest kwadratem.
Lemat~\ref{aa:lem:exceptional-presentation} wraz z szeregiem Hilberta
$2I$ kończy dowód.
\end{proof}

\section{Postacie równań typu E}

\begin{wniosek}[Trzy równania powierzchni]
\label{aa:wn:exceptional-ade}
Przez skalowanie generatorów niezerowymi liczbami
zespolonymi prezentacje z
Twierdzeń~\ref{aa:tw:exceptional-e6-e7}
i~\ref{aa:tw:exceptional-e8} można sprowadzić do
\[
\begin{array}{c|l}
 2T & X^2+Y^3+Z^4=0\\
 2O & X^2+Y^3+YZ^3=0\\
 2I & X^2+Y^3+Z^5=0 .
\end{array}
\]
Są to standardowe postacie równań osobliwości
Kleina typu $E_6,E_7,E_8$.
\end{wniosek}
\begin{proof}
Dla $2T,2I$ współczynniki przy trzech potęgach
są niezerowe i mogą być unormowane do jedynki
przez niezależny wybór odpowiednich pierwiastków.
Dla $2O$ w równaniu $W^2-UV^3+108U^3=0$
przyjmujemy $X=W$, $Y=aU$, $Z=bV$,
gdzie $a^3=108$, $ab^3=-1$.
Daje to drugie równanie. W każdym przypadku
gradient wielomianu definiującego znika tylko
w początku. Dla drugiego równania warunki to
$X=0$, $3Y^2+Z^3=0$ i $3YZ^2=0$.
Ostatnia równość daje $Y=0$ lub $Z=0$,
a poprzednia wymusza oba zera. Dla pierwszego
i trzeciego równania wynika to bezpośrednio
z pochodnych potęg.
\end{proof}

Klasyfikacja i rezolucje osobliwości oraz odpowiedniość McKaya
wymagają dalszych konstrukcji geometrycznych i reprezentacyjnych.
